% 第3章3.4.3迭代算例（verify.py 验算）。弧序 s->a,a->t,s->t；容量 (3,2,4)，费用 (1,1,4)；b_s=3,b_t=-3。
% (a) T={s->t,a->t}，x=(0,0,3)，pi=(4,1,0)，c̄_sa=-2；环上剩余量 3,2,3，步长2。
% (b) a->t 出基进入 U：x=(2,2,1)，费用8，pi=(4,3,0)，c̄_at=-2<=0。
% (c) 若 u_sa=1：步长1，入弧自身到上界，界翻转，x=(1,1,2)，费用10，pi 不变，c̄_sa=-2<=0。
\begin{tikzpicture}[>=Stealth,line width=.8pt,font=\fontsize{8}{10}\selectfont,
  v/.style={circle,draw=FigureStroke,fill=FigureBlueFill,line width=.8pt,minimum size=6.5mm,inner sep=0pt,text=FigureInk},
  rv/.style={v,fill=FigureYellowFill},
  tr/.style={->,draw=FigureBlue,line width=1.4pt,shorten >=.5pt},
  lo/.style={->,draw=FigureMuted,line width=.8pt,dashed,shorten >=.5pt},
  up/.style={->,draw=FigureMuted,line width=.9pt,densely dotted,shorten >=.5pt},
  lab/.style={text=FigureInk,inner sep=1.5pt},
  pot/.style={text=FigureGreen},
  note/.style={text=FigureMuted,font=\FigureNoteFont\fontsize{7.5}{9}\selectfont}]
  \foreach \dx/\k in {0/A,4.9/B,9.8/C}{
  \begin{scope}[xshift=\dx cm]
  \node[v] (s\k) at (0,0) {$s$};
  \node[v] (a\k) at (1.7,1.4) {$a$};
  \node[v] (t\k) at (3.4,0) {$t$};
  \end{scope}}
  % (a)
  \draw[lo,draw=FigureGreen,line width=1.2pt] (sA)--node[lab,above left=-1pt,align=right]{$0/3$\\$+$\,剩$3$} (aA);
  \draw[tr] (aA)--node[lab,above right=-1pt,align=left]{$0/2$\\$+$\,剩$2$} (tA);
  \draw[tr] (sA)--node[lab,below]{$3/4$\quad$-$\,剩$3$} (tA);
  \node[pot,below=1pt] at (sA.south) {$\pi=4$};
  \node[pot,above=1pt] at (aA.north) {$\pi=1$};
  \node[pot,below=1pt] at (tA.south) {$\pi=0$};
  \node[note] at (1.7,-.95) {$\bar c_{s,a}=-2$，$\theta=2$};
  \node at (1.7,-1.5) {(a) 初始树基，费用$12$};
  % (b)
  \draw[tr] (sB)--node[lab,above left=-1pt]{$2/3$} (aB);
  \draw[up,draw=FigureRed,line width=1.1pt] (aB)--node[lab,above right=-1pt]{$2/2$\,$U$} (tB);
  \draw[tr] (sB)--node[lab,below]{$1/4$} (tB);
  \node[pot,below=1pt] at (sB.south) {$\pi=4$};
  \node[pot,above=1pt] at (aB.north) {$\pi=3$};
  \node[pot,below=1pt] at (tB.south) {$\pi=0$};
  \node[note] at (6.6,-.95) {$\bar c_{a,t}=-2\leq0$};
  \node at (6.6,-1.5) {(b) 树弧交换，费用$8$};
  % (c)
  \draw[up,draw=FigureGreen,line width=1.1pt] (sC)--node[lab,above left=-1pt]{$1/1$\,$U$} (aC);
  \draw[tr] (aC)--node[lab,above right=-1pt]{$1/2$} (tC);
  \draw[tr] (sC)--node[lab,below]{$2/4$} (tC);
  \node[pot,below=1pt] at (sC.south) {$\pi=4$};
  \node[pot,above=1pt] at (aC.north) {$\pi=1$};
  \node[pot,below=1pt] at (tC.south) {$\pi=0$};
  \node[note] at (11.5,-.95) {$\bar c_{s,a}=-2\leq0$};
  \node at (11.5,-1.5) {(c) $u_{s,a}=1$时界翻转，费用$10$};
  \node[note] at (6.6,-2.1) {粗蓝线：树弧；虚线：$L$；点线：$U$；绿色：入弧；红色：出弧；绿色数字为节点势；弧旁为$x/u$};
\end{tikzpicture}
